%%%%%%%%%%%%%%%%%%%%%%%%%%%%%%%%%%%%%%%%%%%%%%%%%%%%%%%%%%%
% Matrix Class
%%%%%%%%%%%%%%%%%%%%%%%%%%%%%%%%%%%%%%%%%%%%%%%%%%%%%%%%%%%
%%%%%%%%%%%%%%%%%%%
% Matrix Creation %
%%%%%%%%%%%%%%%%%%%
% Create a new matrix from the given values.
% Equation: \(#1 = \mathbf{M} \in \R^{m\times n} = \begin{bmatrix} #2_{11} & \dots & #2_{1n}\\ \vdots & \ddots & \vdots\\ #2_{m1} & \dots & #2_{mn}\end{bmatrix}\)
% Example: \textbackslash newMat\{myNewMatrix\}\{\{\{0.5,2,3\}\{4,5,6\}\{7,8,9\}\}\}
% Example: \textbackslash myNewMatrix \(\mapsto{} \begin{pmatrix}0.5&2&3\\4&5&6\\7&8&9\end{pmatrix}\)
% Example: \textbackslash newMat\{myNewMatrix\}\{\{\{0.5,2,3\}\}\}
% Example: \textbackslash myNewMatrix \(\mapsto{} \begin{pmatrix}0.5&2&3\end{pmatrix}\)
% Example: \textbackslash newMat\{myNewMatrix\}\{\{\{0.5\}\{4\}\{7\}\}\}
% Example: \textbackslash myNewMatrix \(\mapsto{} \begin{pmatrix}0.5\\4\\7\end{pmatrix}\)
% #1 out, var name to save the result (matrix) to
% #2 scalars, matrix values that are assigned to #1
\newcommand\newMat[2]{
    \xintFor* ##1 in #2 \do {%
        \newVec{matRow}{##1}
        \xintifForFirst{%
            \newMatFromRow{#1}{\matRow}
        }{%
            \appendRowToMat{#1}{\csname #1\endcsname}{\matRow}
        }
    }
}

% TODO matRow should be #2?
% Create a new matrix from the given row.
% #1 out, var name to save the result (matrix) to
% #2 row, matrix row that is assigned to #1
\newcommand\newMatFromRow[2]{
    \newScalar{tmpWidth}{\expandafter\xintLength\expandafter{#2}}%
    \FPifeq{\tmpWidth}{1}
        \expandafter\edef\csname #1\endcsname{{\xintCSVtoList{\matRow}}}
    \else
        \expandafter\edef\csname #1\endcsname{\xintCSVtoList{\matRow}}
    \fi
}

% Create a new matrix from the given row vectors.
% Example: \textbackslash newVec\{myRowVecA\}\{1, 0, 0\}
% Example: \textbackslash newVec\{myRowVecB\}\{0, 2, 0\}
% Example: \textbackslash newVec\{myRowVecC\}\{0, 0, 3\}
% Example: \textbackslash newMatFromRowVecs\{myNewMatrix\}\{\textbackslash myRowVecA, \textbackslash myRowVecB, \textbackslash myRowVecC\}
% Example: \textbackslash myNewMatrix \(\mapsto{} \begin{pmatrix}1&0&0\\0&2&0\\0&0&3\end{pmatrix}\)
% #1 out, var name to save the result (matrix) to
% #2 rows, the rows that are used to construct the matrix
\newcommand\newMatFromRowVecs[2]{%
    % TODO remove this?
    \FPset\numParam{\expandafter\xintLength{#2}}%
    \FPsub\numParam\numParam{1}
    \FPdiv\numParam\numParam{2}
    \FPadd\numParam\numParam{1}
    % In case only one item/row was given the
    % matrix-datastructure must be adjusted
    \FPifeq{\numParam}{1}
        \expandafter\edef\csname #1\endcsname{\xintCSVtoList{#2}}%
    \else
        \expandafter\edef\csname #1\endcsname{\xintCSVtoList{#2}}%
    \fi
}

% Create a new matrix from the given column vectors.
% #1 out, var name to save the result (matrix) to
% #2 columns, the columns that are used to construct the matrix
\newcommand\newMatFromColVecs[2]{%
    \edef\tmpVecList{\xintCSVtoList{#2}}%
    \copyVec[tmpHeightVec]{{\xintNthElt{1}{\tmpVecList}}}
    \getVecHeight[numRows]{\tmpHeightVec}
    \edef\numColumns{\expandafter\xintLength\expandafter{\tmpVecList}}
    \edef\colSequence{\xintSeq[+1]{+1}{\numColumns}}
    \edef\rowSequence{\xintSeq[+1]{+1}{\numRows}}
    \newEmptyMat{#1}
    \xintFor* ##1 in \colSequence \do {%
        \xintFor* ##2 in \rowSequence \do {%
            \newScalar{tmpValue}{\xintNthElt{##1}{\xintNthElt{##2}{\tmpVecList}}}
            \xintifForFirst{%
                \newVec{matRow}{\tmpValue}
            }{%
                \appendToVec{matRow}{\matRow}{\tmpValue}
            }
        }
        \appendRowToMat{#1}{\csname #1\endcsname}{\matRow}
    }
}

% Create a new empty matrix.
% Obacht: It is not necessary to call this to create a new matrix before using any of the arithmetic commands.
% Equation: \(#1 = \mathbf{M} = \emptyset\)
% Example: \textbackslash newEmptyMat\{myNewMatrix\}
% Example: \textbackslash myNewMatrix \(\mapsto{} \emptyset\)
% #1 out, var name to save the result (matrix) to
\newcommand\newEmptyMat[1]{
    \expandafter\edef\csname #1\endcsname{}%
}

% TODO adjust to variable length
% Create a new diagonal matrix from the given vector.
% Equation: \(#1 = \mathbf{M} \in \R^{3\times 3} = \begin{bmatrix} #2_1 & 0 & 0\\ 0 & #2_2 & 0\\ 0 & 0 & #2_3\end{bmatrix}\)
% Example: \textbackslash newDiagMat\{myNewMatrix\}\{0.5,2,3\}
% Example: \textbackslash myNewMatrix \(\mapsto{} \begin{pmatrix}0.5&0&0\\0&2&0\\0&0&3\end{pmatrix}\)
% #1 out, var name to save the result (matrix) to
% #2 scalars, the values that define the diagonal of the matrix
\newcommand\newDiagMat[2]{
    \if\noexpand#1\relax
        \PackageError{glmatrix package}{argument mismatch}{Argument 1 of newDiagMat seems to be a command but must be a string. If it is not a command you might have used two identical letter at the beginning of the argument, please change that.}
    \else
        \edef\tmpVector{\xintCSVtoList{#2}}%
        \newVec{rowA}{\xintNthElt{1}{\tmpVector}, 0, 0}
        \newVec{rowB}{0, \xintNthElt{2}{\tmpVector}, 0}
        \newVec{rowC}{0, 0, \xintNthElt{3}{\tmpVector}}
        \newMatFromRowVecs{#1}{\rowA, \rowB, \rowC}
    \fi
}

% Append the given row to the given matrix.
% Equation: \(#1 = \begin{bmatrix}#2 \\#3\end{bmatrix} \in \R^{m+1\times n}\ \vert\ #2 \in \R^{m\times n}, #3 \in \R^{1\times n}\)
% #1 out, var name to save the result (matrix) to
% #2 matrix, the matrix to append the vector to
% #3 vector r, the row vector that should be appended to the matrix
\newcommand\appendRowToMat[3]{
    \newScalar{tmpWidth}{\expandafter\xintLength\expandafter{#3}}%
    \FPifeq{\tmpWidth}{1}
        \expandafter\edef\csname #1\endcsname{#2 {\xintCSVtoList{#3}}}%
    \else
        \expandafter\edef\csname #1\endcsname{#2 \xintCSVtoList{#3}}%
    \fi
}

%%%%%%%%%%%%%%%%%%%%%%%%%%%
% Transformation Matrices %
%%%%%%%%%%%%%%%%%%%%%%%%%%%
% Create translation matrix from the given vector.
% Equation: \(#1 = \mathbf{T} \in \R^{4\times 4} = \begin{bmatrix}1 & 0 & 0 & #2_x \\0 & 1 & 0 & #2_y \\0 & 0 & 1 & #2_z \\0 & 0 & 0 & 1\end{bmatrix}\)
% Example: \textbackslash newVec\{vecTest\}\{5, 2, 3\}
% Example: \textbackslash newTransMat\{myNewMatrix\}\{\textbackslash vecTest\}
% Example: \textbackslash myNewMatrix \(\mapsto{} \begin{pmatrix}1&0&0&5\\0&1&0&2\\0&0&1&3\\0&0&0&1\end{pmatrix}\)
% #1 out, var name to save the result (matrix) to
% #2 vector t, the translation vector to construct the matrix from
\newcommand\newTransMat[2]{%
    \edef\vecLength{\expandafter\xintLength\expandafter{#2 1}}
    \edef\innersequence{\xintSeq[+1]{+1}{\expandafter\xintLength\expandafter{#2}+1}}
    \xintFor* ##1 in \innersequence \do {%
        \xintifForFirst{%
            \xintFor* ##2 in \innersequence \do {%
                \newScalar{tmp}{0}
                \FPifeq{##2}{\vecLength}
                    \setScalar{tmp}{\xintNthElt{##1}{#2}}
                \else
                \fi
                \FPifeq{##1}{##2}
                    \setScalar{tmp}{1}
                \else
                \fi
                \xintifForFirst{%
                    \newVec{matRow}{\tmp}
                }{%
                    \appendToVec{matRow}{\matRow}{\tmp}
                }
            }
            \newMatFromRow{#1}{\matRow}
        }{%
            \xintFor* ##2 in \innersequence \do {%
                \newScalar{tmp}{0}
                \FPifeq{##2}{\vecLength}
                    \setScalar{tmp}{\xintNthElt{##1}{#2}}
                \else
                \fi
                \FPifeq{##1}{##2}
                    \setScalar{tmp}{1}
                \else
                \fi
                \xintifForFirst{%
                    \newVec{matRow}{\tmp}
                }{%
                    \appendToVec{matRow}{\matRow}{\tmp}
                }
            }
            \appendRowToMat{#1}{\csname #1\endcsname}{\matRow}
        }
    }
}

% Create scaling matrix from the given vector.
% Equation: \(#1 = \mathbf{S} \in \R^{4\times 4} = \begin{bmatrix}#2_x & 0 & 0 & 0 \\0 & #2_y & 0 & 0 \\0 & 0 & #2_z & 0 \\0 & 0 & 0 & 1\end{bmatrix}\)
% Example: \textbackslash newVec\{vecTest\}\{5, 2, 3\}
% Example: \textbackslash newScaleMat\{myNewMatrix\}\{\textbackslash vecTest\}
% Example: \textbackslash myNewMatrix \(\mapsto{} \begin{pmatrix}5&0&0&0\\0&2&0&0\\0&0&3&0\\0&0&0&1\end{pmatrix}\)
% #1 out, var name to save the result (matrix) to
% #2 vector s, the scaling vector to construct the matrix from
\newcommand\newScaleMat[2]{%
    \edef\vecLength{\expandafter\xintLength\expandafter{#2 1}}
    \edef\innersequence{\xintSeq[+1]{+1}{\expandafter\xintLength\expandafter{#2}+1}}
    \xintFor* ##1 in \innersequence \do {%
        \xintifForFirst{%
            \xintFor* ##2 in \innersequence \do {%
                \FPiflt{##2}{\vecLength}
                    \setScalar{tmp}{\xintNthElt{##1}{#2}}
                \else
                    \setScalar{tmp}{1}
                \fi
                \FPifeq{##1}{##2}
                \else
                    \setScalar{tmp}{0}
                \fi
                \xintifForFirst{%
                    \newVec{matRow}{\tmp}
                }{%
                    \appendToVec{matRow}{\matRow}{\tmp}
                }
            }
            \newMatFromRow{#1}{\matRow}
        }{%
            \xintFor* ##2 in \innersequence \do {%
                \FPiflt{##2}{\vecLength}
                    \setScalar{tmp}{\xintNthElt{##1}{#2}}
                \else
                    \setScalar{tmp}{1}
                \fi
                \FPifeq{##1}{##2}
                \else
                    \setScalar{tmp}{0}
                \fi
                \xintifForFirst{%
                    \newVec{matRow}{\tmp}
                }{%
                    \appendToVec{matRow}{\matRow}{\tmp}
                }
            }
            \appendRowToMat{#1}{\csname #1\endcsname}{\matRow}
        }
    }
}

% Create rotation matrix from the given scalar (2D).
% Equation: \(#1 = \mathbf{R} \in \R^{3\times 3}\)
% Example: \textbackslash newScalarPi\{scalarTest\}\{0.5\}
% Example: \textbackslash newRotMatTwoD[myNewMatrix]\{\textbackslash scalarTest\}
% Example: \textbackslash myNewMatrix \(\mapsto{} \begin{pmatrix}0&-1&0\\1&0&0\\0&0&1\end{pmatrix}\)
% #1 out, var name to save the result (matrix) to
% #2 scalar, the rotation value
\newcommand\newRotMatTwoD[2][rotMatResult]{%
    \FPcos\tmpCos{#2}
    \FPsin\tmpSin{#2}
    \newVec{rowB}{\tmpSin, \tmpCos, 0}
    \FPmul\tmpSin\tmpSin{-1}
    \newVec{rowA}{\tmpCos, \tmpSin, 0}
    \newVec{rowC}{0, 0, 1}
    \newMatFromRowVecs{#1}{\rowA, \rowB, \rowC}
}

% Create a new lookAt matrix.
% Equation: \(#1 = \mathbf{L} \in \R^{4\times 4} = \begin{bmatrix}\mathbf{r}_x & \mathbf{r}_y & \mathbf{r}_z & -\mathbf{r} \cdot \mathbf{p} \\\mathbf{u}_x & \mathbf{u}_y & \mathbf{u}_z & -\mathbf{u} \cdot \mathbf{p} \\-\mathbf{d}_x & -\mathbf{d}_y & -\mathbf{d}_z & \mathbf{d} \cdot \mathbf{p} \\0 & 0 & 0 & 1\end{bmatrix}\)
% Example: \textbackslash newVec\{right\}\{5, 2, 3\}
% Example: \textbackslash newVec\{up\}\{0.5, 1.75, 2\}
% Example: \textbackslash newVec\{dir\}\{0.5, 2, -3\}
% Example: \textbackslash newVec\{pos\}\{5, 2, 3\}
% Example: \textbackslash newLookAtMat[myNewMatrix]\{\textbackslash right\}\{\textbackslash up\}\{\textbackslash dir\}\{\textbackslash pos\}
% Example: \textbackslash myNewMatrix \(\mapsto{} \begin{pmatrix}5&2&3&-38\\0.5&1.75&2&-12\\-0.5&-2&3&-2.5\\0&0&0&1\end{pmatrix}\)
% #1 out, var name to save the result (matrix) to
% #2 right, the right vector
% #3 up, the up vector
% #4 dir, the direction vector
% #5 pos, the camera position vector
\newcommand\newLookAtMat[5][lookAtMatResult]{%
    \dotVec[scalarRC]{#2}{#5}
    \mulScalar[scalarRC]{\scalarRC}{-1}
    \copyVec[tmpVecA]{#2}
    \appendToVec{tmpVecA}{\tmpVecA}{\scalarRC}
    % \printVec[1]{\tmpVecA}
    %
    \dotVec[scalarUC]{#3}{#5}
    \mulScalar[scalarUC]{\scalarUC}{-1}
    \copyVec[tmpVecB]{#3}
    \appendToVec{tmpVecB}{\tmpVecB}{\scalarUC}
    %
    \dotVec[scalarDC]{#4}{#5}
    \copyVec[tmpVecC]{#4}
    \negateVec[tmpVecC]{\tmpVecC}
    \appendToVec{tmpVecC}{\tmpVecC}{\scalarDC}
    %
    \newVec{tmpVecD}{0, 0, 0, 1}
    \newMatFromRowVecs{#1}{\tmpVecA, \tmpVecB, \tmpVecC, \tmpVecD}
}

% Create a new projection matrix.
% Equation: \(#1 = \mathbf{P} \in \R^{4\times 4}\)
% Example: \textbackslash newProjMat\{\textbackslash myNewMatrix\}
% Example: \textbackslash myNewMatrix \(\mapsto{} \begin{pmatrix}1&0&0&0\\0&1&0&0\\0&0&1&0\\0&0&-1&0\end{pmatrix}\)
% #1 out, var name to save the result (matrix) to
\newcommand\newProjMat[1]{%
    \newVec{tmpVecA}{1, 0, 0, 0}
    \newVec{tmpVecB}{0, 1, 0, 0}
    \newVec{tmpVecC}{0, 0, 1, 0}
    \newVec{tmpVecD}{0, 0, -1, 0}
    \newMatFromRowVecs{#1}{\tmpVecA, \tmpVecB, \tmpVecC, \tmpVecD}
}

% Create a new frustum matrix.
% Equation: \(#1 = \mathbf{F} \in \R^{4\times 4}\)
% Example: \textbackslash newScalar\{fovx\}\{90\}
% Example: \textbackslash newScalar\{fovy\}\{90\}
% Example: \textbackslash newScalar\{near\}\{1\}
% Example: \textbackslash newScalar\{far\}\{6\}
% Example: \textbackslash newFrustumMat[myNewMatrix]\{\textbackslash fovx\}\{\textbackslash fovy\}\{\textbackslash near\}\{\textbackslash far\}
% Example: \textbackslash myNewMatrix \(\mapsto{} \begin{pmatrix}0.6174&0&0&0\\0&06174&0&0\\0&0&-1.4&-2.4\\0&0&-1&0\end{pmatrix}\)
% #1 out, var name to save the result (matrix) to
% #2 fovx, the fovx scalar
% #3 fovy, the fovy scalar
% #4 near, the position of the near plane (scalar)
% #5 far, the psition of the far lane (scalar)
\newcommand\newFrustumMat[5][frustMatResult]{%
    \FPdiv\tmp{#2}{2}
    \FPtan\frustumRight\tmp
    \FPmul\frustumRight{#4}\frustumRight
    \FPmul\frustumLeft\frustumRight{-1}
    \FPdiv\tmp{#3}{2}
    \FPtan\frustumTop\tmp
    \FPmul\frustumTop{#4}\frustumTop
    \FPmul\frustumBottom\frustumTop{-1}
    %
    \FPmul\cellAA{2}{#4} % AA
    \FPsub\tmp\frustumRight\frustumLeft
    \FPdiv\cellAA\cellAA\tmp
    \FPadd\cellAC\frustumRight\frustumLeft % AC
    \FPdiv\cellAC\cellAC\tmp
    \FPmul\cellBB{2}{#4} % BB
    \FPsub\tmp\frustumTop\frustumBottom
    \FPdiv\cellBB\cellBB\tmp
    \FPadd\cellBC\frustumTop\frustumBottom % BC
    \FPdiv\cellBC\cellBC\tmp
    \FPmul\cellCC{#5}{-1} % CC
    \FPsub\cellCC\cellCC{#4}
    \FPsub\tmp{#5}{#4}
    \FPdiv\cellCC\cellCC\tmp
    \FPmul\cellCD{#5}{#4} % CD
    \FPmul\cellCD{-2}\cellCD
    \FPdiv\cellCD\cellCD\tmp
    %
    \newVec{rowA}{\cellAA, 0, \cellAC, 0}
    \newVec{rowB}{0, \cellBB, \cellBC, 0}
    \newVec{rowC}{0, 0, \cellCC, \cellCD}
    \newVec{rowD}{0, 0, -1, 0}
    \newMatFromRowVecs{#1}{\rowA, \rowB, \rowC, \rowD}
}

% clone
% copy
% create

% fromMat2d
% fromMat4
% fromQuat
% fromQuat2
% fromRotation
% fromRotationTranslation
% fromRotationTranslationScale
% fromRotationTranslationScaleOrigin
% fromXRotation
% fromYRotation
% fromZRotation

% frustum
% lookAt
% ortho
% orthoNO
% orthoZO
% perspective
% perspectiveFromFieldOfView
% perspectiveNO
% perspectiveZO

% normalFromMat4

%%%%%%%%%%%%%%%%%%%
% Matrix Printing %
%%%%%%%%%%%%%%%%%%%
% Print the values of this matrix variable.
% Equation: \(\begin{bmatrix} a_{11} & \dots & a_{1n}\\ \vdots & \ddots & \vdots\\ a_{m1} & \dots & a_{mn}\end{bmatrix}\)
% Example: \textbackslash newMat\{myNewMatrix\}\{\{\{\(a_{11},\dots,a_{1n}\)\}\{\(a_{21},\dots,a_{2n}\)\}\{\(a_{m1},\dots,a_{mn}\)\}\}\}
% Example: \textbackslash printMat\{\textbackslash myNewMatrix\} \(\mapsto{} \begin{bmatrix} a_{11} & \dots & a_{1n}\\ \vdots & \ddots & \vdots\\ a_{m1} & \dots & a_{mn}\end{bmatrix}\)
% #1 fraction, whether or not the value should be printed as a fraction (val=1)
% #2 matrix, the matrix that should be printed
\newcommand\printMat[2][0]{%
    \ensuremath{%
        \begin{bmatrix}
            \xintFor* ##1 in #2 \do {%
                \xintFor* ##2 in {##1} \do {%
                    \xintifForLast{%
                        \xintifForFirst{}{&}
                        \FPifeq{#1}{0}
                            \roundScalar{##2}
                            \scalarRoundResult
                        \else
                            \numToFractionA{##2}
                        \fi
                        \\
                    }{%
                        \xintifForFirst{}{&}
                        \FPifeq{#1}{0}
                            \roundScalar{##2}
                            \scalarRoundResult
                        \else
                            \numToFractionA{##2}
                        \fi
                    }
                }
            }
        \end{bmatrix}
    }
}

% Print only the diagonal values of this matrix variable.
% Example: \textbackslash newMat\{myNewMatrix\}\{\{\{\(a_{11},\dots,a_{1n}\)\}\{\(a_{21},\dots,a_{2n}\)\}\{\(a_{m1},\dots,a_{mn}\)\}\}\}
% Example: \textbackslash printMatDiag\{\textbackslash myNewMatrix\} \(\mapsto{} \text{diag}(a_{11}, a_{22}, a_{mm})\)
% #1 fraction, whether or not the value should be printed as a fraction (val=1)
% #2 matrix, the matrix that should be printed
% TODO variable length
% TODO two params
\newcommand\printMatDiag[2][0]{%
    \ensuremath{%
        \text{diag}\left(
        \xintNthElt{1}{\xintNthElt{1}{#2}},
        \xintNthElt{2}{\xintNthElt{2}{#2}},
        \xintNthElt{3}{\xintNthElt{3}{#2}}
        \right)
    }
}

% Get the height of the given matrix.
% Equation: \(#1 = rows(\mathbf{M})\ \vert\ \mathbf{M} \in \R^{m\times n}\)
% Example: \textbackslash newMat\{myNewMatrix\}\{\{\{\(a_{11},\dots,a_{1n}\)\}\{\(a_{21},\dots,a_{2n}\)\}\{\(a_{m1},\dots,a_{mn}\)\}\}\}
% Example: \textbackslash getMatHeight\{myNewMatrix\}
% Example: \textbackslash matHeightResult \(\mapsto{} m\)
% #1 out, var name to save the result (scalar) to
% #2 matrix, the matrix from which to get the height
\newcommand\getMatHeight[2][matHeightResult]{%
    \newScalar{#1}{\expandafter\xintLength\expandafter{#2}}%
}

% Get the width of the given matrix.
% Equation: \(#1 = cols(\mathbf{M})\ \vert\ \mathbf{M} \in \R^{m\times n}\)
% Example: \textbackslash newMat\{myNewMatrix\}\{\{\{\(a_{11},\dots,a_{1n}\)\}\{\(a_{21},\dots,a_{2n}\)\}\{\(a_{m1},\dots,a_{mn}\)\}\}\}
% Example: \textbackslash getMatWidth\{myNewMatrix\}
% Example: \textbackslash matWidthResult \(\mapsto{} n\)
% #1 out, var name to save the result (scalar) to
% #2 matrix, the matrix from which to get the width
\newcommand\getMatWidth[2][matWidthResult]{%
    \edef\rowElements{\xintNthElt{1}{#2}}%
    \newScalar{#1}{\expandafter\xintLength\expandafter{\rowElements}}%
}

%%%%%%%%%%%%%%%%%%%%%%%
% Matrix Modification %
%%%%%%%%%%%%%%%%%%%%%%%
% Transpose a given Matrix. The parameter can also be a vector,
% but the result will be a matrix.
% Equation: \(#1 = \trans{\mathbf{M}}\)
% Example: \textbackslash newMat\{myNewMatrix\}\{\{\{0.5,2,3\}\{4,5,6\}\{7,8,9\}\}\}
% Example: \textbackslash transposeMat\{\textbackslash myNewMatrix\} \(\mapsto{} \begin{bmatrix}0.5&4&7\\2&5&8\\3&6&9\end{bmatrix}\)
% #1 out, var name to save the result (matrix) to
% #2 matrix, the matrix to transpose
\newcommand\transposeMat[2][transposeMatResult]{%
    \getMatWidth{#2}
    \getMatHeight{#2}
    \edef\widthSequence{\xintSeq[+1]{+1}{\matWidthResult}}
    \edef\heightSequence{\xintSeq[+1]{+1}{\matHeightResult}}
    \newEmptyMat{#1}
    \xintFor* ##1 in \widthSequence \do {%
        \xintFor* ##2 in \heightSequence \do {%
            \FPset\tmpValue{\xintNthElt{##1}{\xintNthElt{##2}{#2}}}
            \xintifForFirst{%
                \newVec{rowVec}{\tmpValue}
            }{%
                \appendToVec{rowVec}{\rowVec}{\tmpValue}
            }
        }
        \appendRowToMat{#1}{\csname #1\endcsname}{\rowVec}
    }
}

% translate
% rotate
% rotateX
% rotateY
% rotateZ
% scale
% set
% projection
% targetTo

%%%%%%%%%%%%%%%
% Matrix Math %
%%%%%%%%%%%%%%%
% Add two or more matrices.
% Equation: \(#1 = \mathbf{X} = #2 + \mathbf{B}\ \vert\ \mathbf{M}, #2, \mathbf{B} \in \R^{m\times n}\)
% Example: \textbackslash newMat\{myNewMatrixA\}\{\{\{0.5,2,3\}\{4,5,6\}\{7,8,9\}\}\}
% Example: \textbackslash newMat\{myNewMatrixB\}\{\{\{0.5,2,3\}\{4,5,6\}\{7,8,9\}\}\}
% Example: \textbackslash addMat\{\textbackslash myNewMatrixA\}\{\textbackslash myNewMatrixB\}
% Example: \textbackslash addMatResult \(\mapsto{} \begin{bmatrix}1&4&6\\8&10&12\\14&16&18\end{bmatrix}\)
% Example: \textbackslash addMat[resultName]\{\textbackslash myNewMatrixA\}\{\textbackslash myNewMatrixB\}
% Example: \textbackslash resultName \(\mapsto{} \begin{bmatrix}1&4&6\\8&10&12\\14&16&18\end{bmatrix}\)
% #1 out, var name to save the result (matrix) to
% #2 matrix, the matrix to add to
% #3 matrices, the matrizes to add
\newcommand\addMat[3][addMatResult]{%
    \if\noexpand#1\relax
        \PackageError{glmatrix package}{argument mismatch}{Argument 1 of addMat seems to be a command but must be a string. If it is not a command you might have used two identical letter at the beginning of the argument, please change that.}
    \else
    \fi
    \if\noexpand#2\relax
    \else
        \PackageError{glmatrix package}{argument mismatch}{argument 2 of addMat does not seem to be a command}
    \fi
    % TODO this does not work if it is a list
    % \if\noexpand#3\relax
    % \else
    %     \PackageError{glmatrix package}{argument mismatch}{argument 3 of addMat does not seem to be a command}
    % \fi
    \def\matlist{\xintCSVtoList{#3}}
    \edef\elemsRow{\xintNthElt{1}{\xintNthElt{1}{\matlist}}}
    \edef\widthsequence{\xintSeq[+1]{+1}{\expandafter\xintLength\expandafter{\elemsRow}}}%
    \edef\elemsMat{\xintNthElt{1}{\matlist}}
    \edef\heightsequence{\xintSeq[+1]{+1}{\expandafter\xintLength\expandafter{\elemsMat}}}%
    %
    \newEmptyMat{#1}
    \xintFor* ##2 in \heightsequence \do {%
        \xintFor* ##3 in \widthsequence \do {%
            \FPset\tmpValue{\xintNthElt{##3}{\xintNthElt{##2}{#2}}}
            \xintifForFirst{%
                \xintFor* ##1 in \matlist \do {%
                    \FPadd\tmpValue\tmpValue{\xintNthElt{##3}{\xintNthElt{##2}{##1}}}
                }
                \newVec{rowVec}{\tmpValue}
            }{%
                % \xintifForLast{
                % }{
                    \xintFor* ##1 in \matlist \do {%
                        \FPadd\tmpValue\tmpValue{\xintNthElt{##3}{\xintNthElt{##2}{##1}}}
                    }
                    \appendToVec{rowVec}{\rowVec}{\tmpValue}
                % }
            }
        }
        \appendRowToMat{#1}{\csname #1\endcsname}{\rowVec}
    }
}

% Subtract two or more matrices.
% Equation: \(#1 = \mathbf{X} = #2 - \mathbf{B}\ \vert\ \mathbf{X}, #2, \mathbf{B} \in \R^{m\times n}\)
% Example: \textbackslash newMat\{myNewMatrixA\}\{\{\{0.5,2,3\}\{4,5,6\}\{7,8,9\}\}\}
% Example: \textbackslash newMat\{myNewMatrixB\}\{\{\{0.25,1,2\}\{3,2,1\}\{3,3,3\}\}\}
% Example: \textbackslash subMat[subMatResult]\{\textbackslash myNewMatrixA\}\{\textbackslash myNewMatrixB\}
% Example: \textbackslash subMatResult \(\mapsto{} \begin{bmatrix}0.25&1&1\\1&3&4\\4&5&6\end{bmatrix}\)
% #1 out, var name to save the result (matrix) to
% #2 matrix, the matrix to subtract from
% #3 matrices, the matrizes to subtract
\newcommand\subMat[3][subMatResult]{%
    % \getMatWidth{#2}
    % \edef\widthsequence{\xintSeq[+1]{+1}{\matWidthResult}}%
    \getMatHeight{#2}
    \edef\heightsequence{\xintSeq[+1]{+1}{\matHeightResult}}%
    %
    \def\matlist{\xintCSVtoList{#3}}
    \edef\elemsRow{\xintNthElt{1}{\xintNthElt{1}{\matlist}}}
    \edef\widthsequence{\xintSeq[+1]{+1}{\expandafter\xintLength\expandafter{\elemsRow}}}%
    % \edef\elemsMat{\xintNthElt{1}{\matlist}}
    % \edef\heightsequence{\xintSeq[+1]{+1}{\expandafter\xintLength\expandafter{\elemsMat}}}%
    %
    \newEmptyMat{#1}
    \xintFor* ##2 in \heightsequence \do {%
        \xintFor* ##3 in \widthsequence \do {%
            \FPset\tmpValue{\xintNthElt{##3}{\xintNthElt{##2}{#2}}}
            \xintifForFirst{%
                \xintFor* ##1 in \matlist \do {%
                    \FPsub\tmpValue\tmpValue{\xintNthElt{##3}{\xintNthElt{##2}{##1}}}
                }
                \newVec{rowVec}{\tmpValue}
            }{%
                \xintFor* ##1 in \matlist \do {%
                    \FPsub\tmpValue\tmpValue{\xintNthElt{##3}{\xintNthElt{##2}{##1}}}
                }
                \appendToVec{rowVec}{\rowVec}{\tmpValue}
            }
        }
        \appendRowToMat{#1}{\csname #1\endcsname}{\rowVec}
    }
}

% Multiply the given matrix with a scalar.
% Equation: \(#1 = \mathbf{X} = #2 * #3\ \vert\ \mathbf{X}, #2 \in \R^{m\times n}, #3 \in \R\)
% Example: \textbackslash newMat\{myNewMatrix\}\{\{\{0.5,2,3\}\{4,5,6\}\{7,8,9\}\}\}
% Example: \textbackslash newScalar\{myNewScalar\}\{2\}
% Example: \textbackslash mulMatScalar\{\textbackslash myNewMatrix\}\{\textbackslash myNewScalar\}
% Example: \textbackslash mulMatScalarResult \(\mapsto{} \begin{bmatrix}1&4&6\\8&10&12\\14&16&18\end{bmatrix}\)
% Example: \textbackslash mulMatScalar[resultName]\{\textbackslash myNewMatrix\}\{\textbackslash myNewScalar\}
% Example: \textbackslash resultName \(\mapsto{} \begin{bmatrix}1&4&6\\8&10&12\\14&16&18\end{bmatrix}\)
% #1 out, var name to save the result (matrix) to
% #2 matrix, the matrix to multiply
% #3 scalar, the scalar to multiply with
\newcommand\mulMatScalar[3][mulMatScalarResult]{%
    \if\noexpand#1\relax
        \PackageError{glmatrix package}{argument mismatch}{Argument 1 of mulMatScalar seems to be a command but must be a string. If it is not a command you might have used two identical letter at the beginning of the argument, please change that.}
    \else
    \fi
    \if\noexpand#2\relax
    \else
        \PackageError{glmatrix package}{argument mismatch}{argument 2 of mulMatScalar does not seem to be a command}
    \fi
    \edef\elemOne{\xintNthElt{1}{#2}}
    \edef\innersequence{\xintSeq[+1]{+1}{\expandafter\xintLength\expandafter{\elemOne}}}%
    %
    \newEmptyMat{#1}
    \xintFor* ##1 in #2 \do {%
        \xintFor* ##2 in \innersequence \do {%
            \FPset\elem{\xintNthElt{##2}{##1}}
            \FPmul\elem\elem{#3}
            \xintifForFirst{%
                \newVec{rowVec}{\elem}
            }{%
                \appendToVec{rowVec}{\rowVec}{\elem}
            }
        }
        \appendRowToMat{#1}{\csname #1\endcsname}{\rowVec}
    }
}

% Multiply the given matrix with a vector (for transformation).
% Equation: \(#1 = #2 \cdot #3\ \vert\ #2 \in \R^{m\times n}, #3 \in \R^n\)
% Error: When matrix and vector have a dimension mismatch
% Example: \textbackslash newVec\{myNewVec\}\{0.5, 1, 2\}
% Example: \textbackslash newMat\{myNewMatrix\}\{\{\{0.5, 1, 2\}\{4, 5, 6\}\{7, 8, 9\}\}\}
% Example: \textbackslash mulMatVec\{\textbackslash myNewMatrix\}\{\textbackslash myNewVec\}
% Example: \textbackslash mulMatVecResult \(\mapsto{} \begin{pmatrix} 5.25\\ 19\\ 29.5\end{pmatrix}\)
% Example: \textbackslash mulMatVec[resultName]\{\textbackslash myNewMatrix\}\{\textbackslash myNewVec\}
% Example: \textbackslash resultName \(\mapsto{} \begin{pmatrix} 5.25\\ 19\\ 29.5\end{pmatrix}\)
% #1 out, var name to save the result (vector) to
% #2 matrix, the matrix to transform the vetor with
% #3 vector, the vector to transform
% TODO if the vector is smaller than the matrix the missing components, should automatically be set to one
\newcommand\mulMatVec[3][mulMatVecResult]{%
    \if\noexpand#1\relax
        \PackageError{glmatrix package}{argument mismatch}{Argument 1 of transformVec seems to be a command but must be a string. If it is not a command you might have used two identical letter at the beginning of the argument, please change that.}
    \else
    \fi
    \if\noexpand#2\relax
    \else
        \PackageError{glmatrix package}{argument mismatch}{argument 2 of transformVec does not seem to be a command}
    \fi
    \if\noexpand#3\relax
    \else
        \PackageError{glmatrix package}{argument mismatch}{argument 3 of transformVec does not seem to be a command}
    \fi
    \getVecHeight{#3}%
    \getMatHeight{#2}%
    \getMatWidth{#2}%
    % Check that matrix and vector can be multiplied
    \FPifeq{\matWidthResult}{\vecHeightResult}
    \else
        \PackageError{glmatrix package}{Dimension mismatch for matrix and vector, could not multiply}{width (\matWidthResult) of matrix #2 and height (\vecHeightResult) of vector #3 do not match}
    \fi
    % Create iterator sequences from height of vector and matrix
    \edef\vecHeightSequence{\xintSeq[+1]{+1}{\vecHeightResult}}%
    \edef\matHeightSequence{\xintSeq[+1]{+1}{\matHeightResult}}%
    % Iterate all rows of the matrix
    \xintFor* ##2 in \matHeightSequence \do {%
        \FPset\cellVal{0}
        % Differentiate between first and other entries to create the new vector variable
        \xintifForFirst{%
            % Now iterate each entry of the current row
            % Use this indec to access the entry in the matrix and in the vector
            \xintFor* ##3 in \vecHeightSequence \do {%
                \FPmul\tmp{%
                    \xintNthElt{##3}{\xintNthElt{##2}{#2}}%
                }{%
                    \xintNthElt{##3}{#3}%
                }
                % Add to storage variable
                \FPadd\cellVal\cellVal\tmp
            }
            \FPclip\cellVal\cellVal
            \newVec{#1}{\cellVal}
        }{
            \xintFor* ##3 in \vecHeightSequence \do {%
                \FPmul\tmp{%
                    \xintNthElt{##3}{\xintNthElt{##2}{#2}}%
                }{%
                    \xintNthElt{##3}{#3}%
                }
                \FPadd\cellVal\cellVal\tmp
            }
            \FPclip\cellVal\cellVal
            \appendToVec{#1}{\csname #1\endcsname}{\cellVal}
        }
    }
    % Test if width and height is correct
    \expandafter\getVecHeight\expandafter{\csname #1\endcsname}%
    \FPifeq{\matHeightResult}{\vecHeightResult}
    \else
        \PackageError{glmatrix package}{Vector height does not match expectation}{height of vector is \vecHeightResult but should be \matHeightResult}
    \fi
}

% Multiply the given matrix with another matrix.
% Error: When the matrizes have a dimension mismatch
% Equation: \(#1 = \mathbf{X} = #2 #3\ \vert\ #2 \in \R^{m\times n}, #3 \in \R^{n\times p}, \mathbf{X} \in \R^{m\times p}\)
% Example: \textbackslash newMat\{myNewMatrixA\}\{\{\{0.5,2,3\}\{4,5,6\}\{7,8,9\}\}\}
% Example: \textbackslash newMat\{myNewMatrixB\}\{\{\{0.25,1,2\}\{3,2,1\}\{3,3,3\}\}\}
% Example: \textbackslash mulMat[mulMatResult]\{\textbackslash myNewMatrixA\}\{\textbackslash myNewMatrixB\}
% Example: \textbackslash mulMatResult \(\mapsto{} \begin{bmatrix}15.125&13.5&12\\34&32&31\\52.75&50&49\end{bmatrix}\)
% #1 out, var name to save the result (matrix) to
% #2 matrix M, the matrix to multiply
% #3 matrix N, the matrix to multiply with
\newcommand\mulMat[3][mulMatResult]{%
    \edef\elemsRow{\xintNthElt{1}{#3}}
    \edef\widthsequence{\xintSeq[+1]{+1}{\expandafter\xintLength\expandafter{\elemsRow}}}%
    %
    % TODO remove one of these
    \edef\elemsMat{#2}
    \edef\heightsequence{\xintSeq[+1]{+1}{\expandafter\xintLength\expandafter{\elemsMat}}}%
    %
    \edef\elemsMat{#3}
    \edef\heightsequenceB{\xintSeq[+1]{+1}{\expandafter\xintLength\expandafter{\elemsMat}}}%
    %
    \edef\elemsRowTmp{\xintNthElt{1}{#2}}
    \FPset\widthMatA{\expandafter\xintLength\expandafter{\elemsRowTmp}}%
    \FPset\heightMatB{\expandafter\xintLength\expandafter{\elemsMat}}%
    % TODO check this
    \FPifeq{\widthMatA}{\heightMatB}
    \else
        \PackageError{glmatrix package}{Dimension mismatch for matrix, could not multiply}{width (\widthMatA) of matrix #2 and height (\heightMatB) of matrix #3 do not match}
    \fi
    %
    \newEmptyMat{#1}
    \xintFor* ##2 in \heightsequence \do {%
        \xintFor* ##3 in \widthsequence \do {%
            \FPset\tmpValue{0}
            \xintifForFirst{%
                \xintFor* ##4 in \heightsequenceB \do {%
                    \FPmul\tmptmp{%
                        \xintNthElt{##4}{\xintNthElt{##2}{#2}}%
                    }{%
                        \xintNthElt{##3}{\xintNthElt{##4}{#3}}%
                    }
                    \FPadd\tmpValue\tmpValue\tmptmp
                }
                \newVec{rowVec}{\tmpValue}
            }{%
                \xintFor* ##4 in \heightsequenceB \do {%
                    \FPmul\tmptmp{%
                        \xintNthElt{##4}{\xintNthElt{##2}{#2}}%
                    }{%
                        \xintNthElt{##3}{\xintNthElt{##4}{#3}}%
                    }
                    \FPadd\tmpValue\tmpValue\tmptmp
                }
                \appendToVec{rowVec}{\rowVec}{\tmpValue}
            }
        }
        \appendRowToMat{#1}{\csname #1\endcsname}{\rowVec}
    }
    % Test if width and height is correct
    \edef\expectedheight{\expandafter\xintLength\expandafter{#2}}%
    \edef\elemsRow{\xintNthElt{1}{#3}}
    \edef\expectedwidth{\expandafter\xintLength\expandafter{\elemsRow}}%
    %
    \edef\elemsRow{\xintNthElt{1}{\csname #1\endcsname}}
    \edef\actualwidth{\expandafter\xintLength\expandafter{\elemsRow}}%
    \edef\elemsMat{\csname #1\endcsname}
    \edef\actualheight{\expandafter\xintLength\expandafter{\elemsMat}}%
    %
    \FPifeq{\expectedwidth}{\actualwidth}
    \else
        \PackageError{glmatrix package}{Matrix width does not match expectation}{width seems to be incorrect}
    \fi
    \FPifeq{\expectedheight}{\actualheight}
    \else
        \PackageError{glmatrix package}{Matrix height does not match expectation}{height seems to be incorrect}
    \fi
}

% multiplyScalarAndAdd

%%%%%%%%%%%%%%%%%%%%%
% Matrix Comparison %
%%%%%%%%%%%%%%%%%%%%%
% equals
% getRotation
% getScaling

% Get the translation component from the given matrix.
% Equation: \(#1 = \mathbf{v} = trans(#2)\ \vert\ #2 \in \R^{m\times n}, \mathbf{v} \in \R^{m-1}\)
% Example: \textbackslash newMat\{myNewMatrix\}\{\{\{0.5,2,3\}\{4,5,6\}\{7,8,9\}\}\}
% Example: \textbackslash getMatTranslation[matTranslationResult]\{\textbackslash myNewMatrix\}
% Example: \textbackslash matTranslationResult \(\mapsto{} \begin{pmatrix}3\\6\\\end{pmatrix}\)
% #1 out, var name to save the result (vector) to
% #2 matrix, the matrix to extract the translation from
\newcommand\getMatTranslation[2][matTranslationResult]{%
    \FPifeq{\expectedheight}{\actualheight}
    \else
        \PackageError{glmatrix package}{Matrix height does not match expectation}{height seems to be incorrect}
    \fi
    \getMatWidth{#2}
    \getMatHeight{#2}
    \edef\sequence{\xintSeq[+1]{+1}{\matHeightResult}}
    \xintFor* ##1 in \sequence \do {%
        \FPset\tmpValue{\xintNthElt{\matWidthResult}{\xintNthElt{##1}{#2}}}
        \xintifForFirst{%
            \newVec{#1}{\tmpValue}
        }{%
            \xintifForLast{}{
                \appendToVec{#1}{\csname #1\endcsname}{\tmpValue}
            }
        }
    }
}

%%%%%%%%%%%%%%%
% Matrix Misc %
%%%%%%%%%%%%%%%
% Get the determinate of the given matrix.
% Equation: \(#1 = det(#2)\ \vert\ #2 \in \R^{m\times n}\)
% Example: \textbackslash newMat\{myNewMatrix\}\{\{\{0.5,2,3\}\{4,5,6\}\{7,8,9\}\}\}
% Example: \textbackslash detMat[detMatResult]\{\textbackslash myNewMatrix\}
% Example: \textbackslash detMatResult \(\mapsto{} 1.5\)
% #1 out, var name to save the result (scalar) to
% #2 matrix, the matrix to calculate the determinate for
\newcommand\detMat[2][detMatResult]{%
    \getMatWidth{#2}
    \getMatHeight{#2}
    \FPifgt{\matWidthResult}{2}
        \PackageError{glmatrix package}{Function not yet implemented}{Determinante can only be calculated for \(2 \times 2\) matrices}
    \else
    \fi
    \FPifgt{\matHeightResult}{2}
        \PackageError{glmatrix package}{Function not yet implemented}{Determinante can only be calculated for \(2 \times 2\) matrices}
    \else
    \fi
    %
    \mulScalar[tmpAD]{\xintNthElt{1}{\xintNthElt{1}{#2}}}{\xintNthElt{2}{\xintNthElt{2}{#2}}}
    \mulScalar[tmpBC]{\xintNthElt{1}{\xintNthElt{2}{#2}}}{\xintNthElt{2}{\xintNthElt{1}{#2}}}
    \subScalar[#1]{\tmpAD}{\tmpBC}
}

% adjoint
% decompose
% determinant
% frob
% identity
% invert

% LDU
% str
